\section{General Inverse Reconstruction}
\label{sec:inverse-reconstruction}

The preceding sections began with a refractive-index profile and derived the
resulting horizon law. The same relation can instead be read in reverse.
Inverse methods have previously been applied to atmospheric-refraction data
\cite{rees1991inversion}; the reconstruction below is independently derived
for the specific grazing-distance relation of the present model. Given a
desired one-sided visibility distance as a function of endpoint height, the
corresponding stratified refractive-index profile can be reconstructed
directly. This inverse formulation separates mathematical existence from
the physical
question of whether the reconstructed profile can occur in an atmosphere.

\subsection{Reconstruction from a prescribed horizon law}

For a monotonic profile with $n(0)=n_0$ and $n(h)>n_0$ for $h>0$, the grazing
ray derived in Sec.~\ref{sec:model} gives the one-sided horizon distance
\begin{equation}
    d(h)
    =
    \int_0^h
    \frac{n_0\,d\eta}
    {\sqrt{n^2(\eta)-n_0^2}}.
    \label{eq:inverse-forward-relation}
\end{equation}
Because the integrand depends only on the local refractive index,
differentiation of Eq.~\eqref{eq:inverse-forward-relation} with respect to
its upper limit yields
\begin{equation}
    d'(h)
    =
    \frac{n_0}
    {\sqrt{n^2(h)-n_0^2}}.
    \label{eq:horizon-law-derivative}
\end{equation}
Solving algebraically for $n(h)$ gives the inverse reconstruction formula
\begin{equation}
    n(h)
    =
    n_0
    \sqrt{
        1+\frac{1}{[d'(h)]^2}
    }.
    \label{eq:general-inverse-reconstruction}
\end{equation}
Thus, within the assumptions of the model, any differentiable and strictly
increasing horizon law with $d'(h)>0$ defines a refractive-index profile. The
overall multiplicative scale $n_0$ does not affect ray trajectories; it fixes
only the absolute optical path length. Once $n_0$ is chosen, the profile is
unique wherever $d'(h)$ is finite and nonzero.

Equation~\eqref{eq:general-inverse-reconstruction} also makes the near-boundary
behaviour transparent. A square-root horizon law,
$d(h)\sim C\sqrt{h}$, has $d'(h)\sim C/(2\sqrt{h})$ and therefore
$n(h)\to n_0$ continuously as $h\to0^+$. Expanding the square root gives
\begin{equation}
    n(h)
    =
    n_0
    \left[
        1+\frac{1}{2[d'(h)]^2}
        +O\!\left(\frac{1}{[d'(h)]^4}\right)
    \right].
    \label{eq:inverse-near-boundary-expansion}
\end{equation}
For $d(h)\simeq\sqrt{2Rh}$, this immediately reproduces
$n(h)\simeq n_0(1+h/R)$ and hence the leading-order exponential profile found
in Sec.~\ref{sec:paraxial}.

\subsection{Exact reconstruction of the spherical horizon law}

We now prescribe the exact spherical arc-distance law from
Eq.~\eqref{eq:exact-spherical-horizon-distance},
\begin{equation}
    d_{\mathrm{s}}(h)
    =
    R\arccos\!\left(\frac{R}{R+h}\right).
\end{equation}
Its derivative is
\begin{equation}
    d_{\mathrm{s}}'(h)
    =
    \frac{R^2}
    {(R+h)\sqrt{h(2R+h)}}.
    \label{eq:spherical-horizon-derivative}
\end{equation}
Substitution into Eq.~\eqref{eq:general-inverse-reconstruction} gives
\begin{equation}
    n_{\mathrm{s}}(h)
    =
    n_0
    \sqrt{
        1+
        \frac{(R+h)^2h(2R+h)}{R^4}
    }.
    \label{eq:exact-spherical-reconstruction-profile}
\end{equation}
By construction, a grazing ray above the flat opaque boundary has exactly the
same one-sided distance--height relation as a tangent
line of sight above a sphere of radius $R$, when both models use the same
numerical height $h$. Therefore
\begin{equation}
    L_{\max}
    =
    d_{\mathrm{s}}(h_{\mathrm{o}})
    +
    d_{\mathrm{s}}(h_{\mathrm{t}})
\end{equation}
is reproduced exactly for arbitrary observer and target heights within the
idealised model.

For $h/R\ll1$, Eq.~\eqref{eq:exact-spherical-reconstruction-profile} expands as
\begin{equation}
    \frac{n_{\mathrm{s}}(h)}{n_0}
    =
    1
    +\frac{h}{R}
    +2\left(\frac{h}{R}\right)^2
    +O\!\left(\frac{h^3}{R^3}\right).
    \label{eq:spherical-profile-small-height}
\end{equation}
The linear term agrees with
$e^{h/R}=1+h/R+\tfrac12(h/R)^2+\cdots$, explaining why the exponential
profile reproduces the spherical law at leading order. The different
quadratic terms encode the small but finite discrepancy quantified in
Sec.~\ref{sec:spherical-comparison}.

\subsection{Scope of the inverse result}

The reconstruction determines the limiting grazing distance and the associated
boundary-occlusion criterion. It does not by itself establish
that every aspect of image formation, intensity, dispersion, turbulence, or
multi-path propagation is identical to that of a spherical surface. Nor does
mathematical reconstructibility imply physical plausibility. In particular,
Eq.~\eqref{eq:exact-spherical-reconstruction-profile} increases with height,
just as the exponential profile does. The next section compares this required
sign and scale with those of a physical atmosphere.
